\begin{definition}
Let $\sigma$ and $a$ be as in \noderef{SubellipticOp}, let $U\subset\mathbb{R}^n$ and let
$u:\mathbb{R}^n\to\mathbb{R}$. We say that $u$ is a \emph{classical sub-solution} of
\texttt{eq:maineq} in $U$, and write $\mathcal{L}u\leq0$ in $U$, if
$$\mathcal{L}u(x)\leq0\qquad\text{at every }x\in U\text{ with }\mathfrak{X}u(x)\neq0 .$$
Here $\mathfrak{X}u$ is the horizontal gradient of \noderef{HGrad} and $\mathcal{L}$ the operator
of \noderef{SubellipticOp}. No condition is imposed at the horizontal critical points of $u$,
where $a$ need not be defined.

This is the sub-solution condition of Definition \texttt{def:classical} of the paper. The paper
additionally takes $U$ open and $u\in C^2(U)$; those are regularity requirements on $U$ and $u$
rather than part of the differential inequality, and the nodes that use this notion carry them as
separate hypotheses.
\end{definition}
